<<<<<<< HEAD
\begin{name}
	{\tenchude}
	{TOÁN 10}
	{LỚP TOÁN THẦY PHÁT}
	{I can accept failure, everyone fails at something. But I can’t accept not trying.}
\end{name}
\Opensolutionfile{ansbook}[ans/ansbookDe8]
\TN
\Opensolutionfile{ans}[ans/ansDe1-TN8]
\begin{ex}%[HSA-đợt 1,Huy Trần- Đăng Tạ]%[0D8N1-5]
Từ các chữ số $1$, $2$, $3$, $4$, $5$ có thể lập được bao nhiêu số tự nhiên có $3$ chữ số đôi một khác nhau?
\choice{\True $60$} {$125$} {$5$} {$10$}
\loigiai{
Lập số tự nhiên có $3$ chữ số đôi một khác nhau từ các chữ số $1$, $2$, $3$, $4$, $5$ có
\[A_3^5 = 60 ~\text {số}.\]
}
\end{ex}

\begin{ex}%[0D8N2-1]%[Dự án đề kiểm tra Toán khối 11 HKII NH23-24-Dot15-PhucChuong]%[THPT Nhu Van Lang - Hai Phong]
Công thức tính số chỉnh hợp chập $k$ của $n$ phần tử là:
\choice
{$A_n^k=\dfrac{n!}{\left(n-k\right)!k!}$}
{$C_n^k=\dfrac{n!}{\left(n-k\right)!}$}
{$C_n^k=\dfrac{n!}{\left(n-k\right)!k!}$}
{\True$A_n^k=\dfrac{n!}{\left(n-k\right)!}$}
\loigiai{
Công thức tính số chỉnh hợp chập $k$ của $n$ phần tử là $A_n^k=\dfrac{n!}{\left(n-k\right)!}$.
}
\end{ex}

\begin{ex}%[0D8N3-4]
Trong khai triển $\left(2a-b\right)^5$, hệ số của số hạng thứ $3$ bằng
\choice
{$-80$}
{\True $80$}
{$-10$}
{$10$}
\loigiai{
$
\begin{aligned}
\text{Ta có}
\left(2a-b\right)^5&=(2a)^5-5(2a)^4b+10(2a)^3b^2-10(2a)^2b^3+5(2a)b^4-b^5 \\
&=32a^5-80a^4b+80a^3b^2-40a^2b^3+10ab^4-b^5.
\end{aligned}
$
}
\end{ex}

\begin{ex}%[H]%[Nguyễn Văn Cường (Cường NV) - Dự án giảng K10 - K11]%[0D6H1-1]
Cho số $x = \dfrac{2}{7}$ và các giá trị gần đúng của $x$ là $0{,}28$; $0{,}29$; $0{,}286$; $0{,}3$. Hãy xác định sai số tuyệt đối trong từng trường hợp và cho biết giá trị gần đúng nào là tốt nhất.
\choice
{$0{,}28$}
{$0{,}29$}
{\True $0{,}286$}
{$0{,}3$}
\loigiai{
Ta có các sai số tuyệt đối là
\[\Delta_a=\left|\dfrac{2}{7}-0{,}28\right|=\dfrac{1}{175}; \Delta_b=\left|\dfrac{2}{7}-0{,}29\right|=\dfrac{3}{700} \Delta_c=\left|\dfrac{2}{7}-0{,}286\right|=\dfrac{1}{3500}; \Delta_d=\left|\dfrac{2}{7}-0{,}3\right|=\dfrac{1}{70}.\]
Vì $\Delta_c < \Delta_b < \Delta_a < \Delta_d$ nên $c=0{,}286$ là số gần đúng tốt nhất.}
\end{ex}

\begin{ex}%[0D6H4-2]
Nhiệt độ trung bình hàng tháng trong một năm được ghi lại trong bảng sau:
\begin{center}
\begin{tabular}{|c|c|c|c|c|c|c|c|c|c|c|c|c|}
\hline
Tháng &1&2&3&4&5&6&7&8&9&10&11&12\\
\hline
Nhiệt độ&16&20&24&28&30&30&29&25&25&20&18&16\\
\hline
\end{tabular}
\end{center}
Tìm khoảng tứ phân vị của bảng số liệu trên.
\choice
{$24,5$}
{$19$}
{$28,5$}
{\True $9,5$}
\loigiai{
Sắp xếp mẫu số liệu theo thứ tự không giảm:
\[16\hspace{1cm}16\hspace{1cm}18\hspace{1cm}20\hspace{1cm}20\hspace{1cm}24\hspace{1cm}25\hspace{1cm}25\hspace{1cm}28\hspace{1cm}29\hspace{1cm}30\hspace{1cm}30\]
Mẫu số liệu gồm $12$ giá trị nên trung vị là $Q_2=\dfrac{24+25}{2}=24,5$.\\
Nửa số liệu bên trái gồm $6$ giá trị. Do đó, $Q_1=\dfrac{18+20}{2}=19$.\\
Nửa số liệu bên phải gồm $6$ giá trị. Do đó, $Q_3=\dfrac{28+29}{2}=28,5$.\\
Vậy khoảng tứ phân vị của bảng số liệu trên là: $\Delta_Q=28,5-19=9,5$.
}
\end{ex}

\begin{ex}%[0D0N1-2]%[Dự án đề kiểm tra Toán 11 GHKII NH23-24- Hieu Hieu Minh Minh]%[THPT Toàn Thắng- Hải Phòng]
Không gian mẫu của phép thử \lq\lq Gieo đồng thời $2$ đồng xu cân đối và đồng chất\rq\rq\, là
\choice
{$\Omega=\left\{S, N, S, N\right\}$}
{\True $\Omega=\left\{SS, NN, SN, NS\right\}$}
{$\Omega=\left\{SS, NN\right\}$}
{$\Omega=\left\{S, N\right\}$}
\loigiai{
Gieo đồng thời $2$ đồng xu cân đối và đồng chất, ta được không gian mẫu là $\Omega=\left\{SS, NN, SN, NS\right\}$.
}
\end{ex}

\begin{ex}%[Dự án 17 - CK2 NH23-24 - TeamTeXHoa - ThinhNgo]%[0D0N1-3]
Gieo đồng tiền $2$ lần. Số phần tử của biến cố để mặt ngửa xuất hiện ít nhất $1$ lần là
\choice
{$6$}
{$5$}
{\True $3$}
{$4$}
\loigiai{
Không gian mẫu $\Omega=\{NN,NS,SN,SS\}$.\\
Gọi $A$ là biến cố mặt ngửa xuất hiện ít nhất $1$ lần. Khi đó $A=\{NN,NS,SN\}$ $\;\Rightarrow\; n(A)=3$.
}
\end{ex}

\begin{ex}%[0H9H3-2]
Trong mặt phẳng $Oxy$, phương trình tổng quát của đường thẳng đi qua hai điểm $A\left(3;-1\right)$ và $B\left(1;5\right)$ là
\choice
{$-x+3y+6=0$}
{$3x-y+6=0$}
{$3x-y+10=0$}
{\True $3x+y-8=0$}
\loigiai{
\begin{align*}
\heva{	& A\left(3;-1\right)\in AB \\
& {{\vec{u}}}_{AB}=\overrightarrow{AB}=\left(-2;6\right)\Rightarrow {{\vec{n}}}_{AB}=\left(3;1\right).}
\end{align*}
Phương trình đường thẳng $AB$ là \[3\left(x-3\right)+1\left(y+1\right)=0\Leftrightarrow 3x+y-8=0.\]
}
\end{ex}

\begin{ex}%[0H9H3-3]
Hai đường thẳng $\left(d_1\right) \colon \heva{&x=-2+5 t \\ &y=2 t}$ và $\left(d_2\right) \colon 4 x+3 y-18=0$ cắt nhau tại điểm có tọa độ
\choice
{\True $(2 ; 3)$}
{$(3 ; 2)$}
{$(1 ; 2)$}
{$(2 ; 1)$}
\loigiai{
Ta có
\[
\left(d_1\right) \colon \heva{&x=-2+5 t \\&y=2 t}
\Rightarrow
\left(d_1\right) \colon 2 x-5 y+4=0
\]
Gọi $ M $ là giao điểm của $ (d_1) $ và $ (d_2) $ nên tọa độ $ M $ là nghiệm của hệ phương trình
\[
\heva{&2 x-5 y+4=0 \\ &4 x+3 y-18=0}
\Leftrightarrow
\heva{&x=2 \\ &y=3.}
\]
Vậy $ M(2;3) $.
}
\end{ex}

\begin{ex}%[0H9H3-5]%[tex hóa đề ck2 - form 2025 - đợt 2- Duong Xuan Loi]
Khoảng cách từ điểm $M(3;-1)$ đến đường thẳng $\Delta\colon\heva{&x=-2+t\\&y=1+2t}$ nằm trong khoảng nào sau đây?
\choice
{$(1;3)$}
{$(3;5)$}
{$(7;9)$}
{\True $(5;7)$}
\loigiai{
Phương trình tổng quát đường thẳng $\Delta$ là $2x-y+5=0$.\\
Khoảng cách từ điểm. $M$. đến đường thẳng $\Delta$ là $\dfrac{\left|2\cdot 3-(-1)+5\right|}{\sqrt{2^2+(-1)^2}}=\dfrac{12\sqrt{5}}{5}\approx 5,4$.}
\end{ex}

\begin{ex}%[0H9N5-1]%[Dự án đề kiểm tra Toán Khối 10 CK2 NH23-24-Dot 16- Khắc Thiên]%[THPT Chuyên Hùng Vương - Phú Thọ]
Trong mặt phẳng $Oxy$, cho đường elip $(E)$ có phương trình  $\dfrac{x^2}{36}+\dfrac{y^2}{9}=1$. Tổng khoảng cách từ mỗi điểm trên elip tới hai tiêu điểm bằng
\choice
{$6$}
{$3$}
{$5$}
{\True $12$}
\loigiai{
Ta có phương trình của chính tắc elip  $(E)$ là $\dfrac{x^2}{36}+\dfrac{y^2}{9}=1$.\\
Suy ra $a=6$.\\
Khi đó, tổng khoảng cách từ mỗi điểm trên elip tới hai tiêu điểm bằng $2a=12$.
}
\end{ex}

\begin{ex}%[0H9N5-5]
Phương trình nào sau đây là phương trình chính tắc của hypebol?
\choice
{$x^2+\dfrac{y^2}{3^2}=1$}
{$\dfrac{x^2}{16}-y^2=-1$}
{$\dfrac{x^2}{25}-\dfrac{y^2}{9}=-1$}
{\True $x^2-\dfrac{y^2}{2}=1$}
\loigiai{
Phương trình
$x^2-\dfrac{y^2}{2}=1 \Leftrightarrow \dfrac{x^2}{1}-\dfrac{y^2}{2}=1$ là phương trình chính tắc của hypebol.
}
\end{ex}
\TNTF
\begin{ex}%[24-25 giảng K10 - K11, Phan Quốc Trí]%[0D8V2-6]
Cho đa giác đều $12$ cạnh. Xét tính đúng sai của các khẳng định sau.
\choiceTF
{\True Đa giác này có $54$ đường chéo}
{\True Có $220$ tam giác được tạo thành từ các đỉnh của đa giác}
{Số hình chữ nhật được tạo thành từ các đỉnh của đa giác là $165$}
{Từ các đỉnh của đa giác, có $120$ tam giác có đúng một cạnh là cạnh của đa giác}
\loigiai{
\begin{itemchoice}
\itemch Đúng. Số đường chéo của đa giác đều $12$ cạnh là:
$\mathrm{C}_{12}^2 - 12=54$ đường chéo.
\itemch Đúng. Số tam giác được tạo thành từ các đỉnh của đa giác là số cách chọn $3$ đỉnh từ $12$ đỉnh:
$\mathrm{C}_{12}^3=220$ tam giác.
\itemch Sai. Đa giác đều này có $6$ đường chéo đi qua tâm. Cứ hai đường chéo qua tâm cho ta một hình chữ nhật. Do đó có $\mathrm{C}_{6}^2 = 15$ hình chữ nhật tạo thành.
\itemch Sai. Trước tiên ta chọn $1$ cạnh trong $12$ cạnh của đa giác nên có $\mathrm{C}_{12}^1$ cách chọn; tiếp theo chọn $1$ đỉnh còn lại trong $8$ đỉnh (trừ $2$ đỉnh tạo nên cạnh đã chọn và $2$ đỉnh liền kề với cạnh đã chọn) nên có $\mathrm{C}_{8}^1$ các chọn. Do đó Số tam giác có một cạnh là cạnh của đa giác là:
$\mathrm{C}_{12}^1\times\mathrm{C}_{8}^1=96$ tam giác.
\end{itemchoice}
}
\end{ex}

\begin{ex}%[0H9H3-2]%[Dự án đề 4 phần]$%[Tex hóa: Lê Thị Thúy Hằng]
Trong mặt phẳng toạ độ $Oxy$, cho hai điểm $A(-2;2)$, $B(3;4)$.
\choiceTF
{Đường thẳng $AB$ có vectơ chỉ phương là $\overrightarrow{AB}=(2;5)$}
{Trung điểm của đoạn thẳng $AB$ có tọa độ là $(1;6)$}
{\True Phương trình tổng quát của đường thẳng $AB$ là $2x-5y+14=0$}
{Phương trình tham số của đường thẳng đi qua $M(-1;1)$ và song song với $AB$ là $\heva{&x=-1+2t\\&y=1+5t}$}
\loigiai{
\begin{itemchoice}
\itemch {\bf Sai}.\\
Đường thẳng $AB$ có vectơ chỉ phương là $\overrightarrow{AB}=(5;2)$.
\itemch {\bf Sai}.\\
Trung điểm của đoạn thẳng $AB$ có tọa độ là $\left(\dfrac{1}{2};3\right)$.
\itemch {\bf Đúng}.\\
Phương trình tổng quát của đường thẳng $AB$ đi qua $A(-2;2)$ và có vectơ pháp tuyến $\overrightarrow{n}=(2;-5)$ là \[2(x+2)-5(y-2)=0 \Leftrightarrow 2x-5y+14=0.\]
\itemch {\bf Sai}.\\
Đường thẳng đi qua $M(-1;1)$, song song với đường thẳng $AB$ nên nhận $\overrightarrow{AB}=(5;2)$ làm vectơ chỉ phương, có phương trình tham số là $\heva{&x=-1+5t\\&y=1+2t.}$
\end{itemchoice}
}
\end{ex}
\TNSA
\begin{ex}%[0D8V1-2]%[Dự án đề kiểm tra Toán 10 CHKII NH23-24- Ngô Tất Thành]%[THPT Chuyên Hùng Vương - Tỉnh Phú Thọ]
Đội thanh niên xung kích của một trường trung học phổ thông có $12$ học sinh trong đó có $9$ học sinh nam và $3$ học sinh nữ. Đoàn trường cần chọn một nhóm $5$ học sinh đi làm nhiệm vụ sao cho phải có $1$ đội trưởng nam, $1$ đội phó nam và có ít nhất $1$ nữ. Hỏi có bao nhiêu cách?
\shortans{$6120$}
\loigiai{
Chọn $1$ nam làm đội trưởng có $\mathrm{C}_9^1$ cách.\\
Chọn $1$ nam làm đội phó trong 8 nam còn lại có $\mathrm{C}_8^1$ cách.\\
Chọn $3$ học sinh sao cho có ít nhất $1$ nữ có $\mathrm{C}_{10}^3-\mathrm{C}_7^3=85$.\\
Số cách chọn $5$ bạn thỏa yêu cầu bài toán là $\mathrm{C}_9^1 \cdot \mathrm{C}_8^1 \cdot 85=6120$ cách.
}
\end{ex}

\begin{ex}%[0D8N3-4]
Tổng các hệ số của các số hạng trong khai triển nhị thức $(x-2y)^{2\,020}$ bằng
\choice{$2\,021$}{$2\,020$}{$-1$}{\True $1$}
\loigiai{
Ta có $(x-2y)^{2\,020}=\sum\limits_{k=0}^{2\,020}\mathrm{C}_{2\,020}^kx^{2\,020-k}(-2y)^k=\sum\limits_{k=0}^{2\,020} a_kx^{2\,020-k}y^k$.\\
Chọn $x=y=1$ ta có $\sum\limits_{k=0}^{2\,020} a_k=(-1)^{2\,020}=1$.
}
\end{ex}

\begin{ex}%[De-chuan-hoa-so-8]%[Duong Xuan Loi]%[0D6H3-2]
Bảng dưới đây thống kê nhiệt độ (đơn vị: ${}^{\circ}{C}$) ở thành phố Hồ Chí Minh ngày $03 / 06 / 2021$ sau một số lần đo
\begin{center}
\begin{tabular}{|l|c|c|c|c|c|c|c|c|}
\hline Giờ đo & $1 $ h
& $4 $ h & $7 $ h & $10 $ h & $13 $ h & $16 $ h & $19 $ h & $22 $ h \\
\hline Nhiệt độ $\left({}^{\circ}{C}\right)$ & $27$ & $26$ & $28$ & $32$ & $34$ & $35$ & $30$ & $28$ \\
\hline
\end{tabular}
\end{center}
Tìm số trung bình của mẫu số liệu (làm tròn kết quả đến hàng phần trăm).
\shortans{$29$}
\loigiai{
Sắp xếp dãy số liệu nhiệt độ theo thứ tự không giảm ta được: $26;27;28;28;30;32;34;35$.\\
Nhiệt độ trung bình $\overline{x}=\dfrac{26+27+28\cdot 2+30+32+34+35}{8}=30$.\\
Số trung vị $M_e=\dfrac{28+30}{2}=29$.
}
\end{ex}

\begin{ex}%[0H9V4-2]
Trong mặt phẳng tọa độ $Oxy$, cho điểm $A(-1;1)$, $B(0;-2)$, $C(0;2)$. Đường tròn ngoại tiếp tam giác $ ABC $ có phương trình là $ (x-a)^2+(y-b)^2=R^2 $. Tính $ S=a+b+R^2 $.
\shortans{$6$}
\loigiai{
Giả sử tâm đường tròn là điểm $I(a;b)$.\\
Ta có \allowdisplaybreaks
\begin{eqnarray*}
& &IA=IB=IC\Leftrightarrow IA^2=IB^2=IC^2\\
&\Leftrightarrow&\heva{& (-1-a)^2+(1-b)^2=(0-a)^2+(-2-b)^2\\ & (0-a)^2+(-2-b)^2=(0-a)^2+(2-b)^2}\\
&\Leftrightarrow& \heva{& a^2+b^2+2a-2b+2=a^2+b^2+4b+4\\ & a^2+b^2+4b+4=a^2+b^2-4b+4}\\
&\Leftrightarrow& \heva{& 2a-2b=4b+2\\ & b=0}\Leftrightarrow\heva{& a=1\\ & b=0.}
\end{eqnarray*}
Đường tròn tâm $I(1;0)$, bán kính $R=IC=\sqrt{a^2+b^2-4b+4}=\sqrt{5}$.\\
Suy ra phương trình đường tròn đi qua ba điểm $A$, $B$, $C$ có dạng $(x-1)^2+y^2=5$.\\
Vậy $a=1$, $b=0$, $R^2=5$, $a+b+R^2=6$.
}
\end{ex}
\TL
\begin{ex}%[0H9H3-2]%[Dự án đề kiểm tra Toán 10 HK2 NH23-24-BCTuan]%[THPT Nguyễn Thị Minh Khai-TPHCM]
Trong mặt phẳng $Oxy$, cho $A(1;2)$, $B(2;3)$, $C(6;6)$. Viết phương trình tổng quát của đường thẳng $\Delta$ đi qua điểm $A$ và song song với $BC$.
\loigiai{
Ta có $\overrightarrow{BC}=(4;3)$ suy ra $\overrightarrow{n}=(3;-4)$ là véc-tơ pháp tuyến của $BC$.\\
Vì $\Delta$ song song với $BC$ nên $\overrightarrow{n}=(3;-4)$ cũng là véc-tơ pháp tuyến của $\Delta$.\\
Ta có phương trình tổng quát của $\Delta$ là $3(x-1)-4(y-2)=0\Leftrightarrow 3x-4y+5=0$.
}
\end{ex}

\begin{ex}%[0H9V4-5]%[Dự án đề kiểm tra Toán 11 GHKI NH23-24- Phan Anh]%[THPT - Tp HCM]
Trong mặt phẳng tọa độ $Oxy$, cho điểm $I(2;-1)$ và đường thẳng $d\colon3x+y+5=0$. Đường tròn $(C)$ có tâm $I$ và đường kính bằng $4\sqrt{5}$. Đường thẳng $\Delta\colon x+by+c=0$, $c>0$ vuông góc với đường thẳng $d$ và cắt $(C)$ tại hai điểm phân biệt $A$, $B$ sao cho tam giác $IAB$ tù và có diện tích bằng $5\sqrt{3}$. Giá trị của $c$ bằng bao nhiêu? (Làm tròn kết quả đến hàng phần chục)
\loigiai{\immini[thm]{Đường tròn $(C)$ có đường kính bằng $4\sqrt{5}$ nên có bán kính là $R=2\sqrt{5}$.\\
Đường thẳng $\Delta$ vuông góc với đường thẳng $d$ nên $\Delta\colon x-3y+c=0$.\\
Tam giác $IAB$ có diện tích là $5\sqrt{3}$ nên
\[S=\dfrac{1}{2}\cdot IA\cdot IB\cdot\sin\widehat{AIB}\Leftrightarrow5\sqrt{3}=\dfrac{R^2\sin\widehat{AIB}}{2}\Leftrightarrow\sin\widehat{AIB}=\dfrac{\sqrt{3}}{2}.\]}
{\begin{tikzpicture}[>=stealth,line join=round,line cap=round,font=\footnotesize,scale=0.3]
\path (0,0) coordinate (I)
($(I)+({5*cos(30)},{5*sin(30)})$) coordinate(A)
($(I)!1!120:(A)$) coordinate(B)
($(A)!0.5!(B)$) coordinate(H)
;
\draw (I)circle(5);
\draw (H)--(I)--(B)--(A)--(I);
\foreach \p / \r in {A/45,I/-90,B/135,H/90}
\fill (\p) circle (1.2pt) node[shift={(\r:3mm)}]{$\p$};
\end{tikzpicture}}
\noindent
Do $\triangle IAB$ tù nên $\widehat{AIB}=120^\circ$, suy ra khoảng cách từ $I$ đến $\Delta$ là $IH=\dfrac{R}{2}=\sqrt{5}$.\\
Từ đó ta có $IH=\mathrm{d}(I,\Delta)=\dfrac{|2+3+c|}{\sqrt{1+9}}=\sqrt{5}\Leftrightarrow\hoac{&c=5\sqrt{2}-5\\&c=-5\sqrt{2}-5.}$\\
Do $c$ dương nên nhận $c=5\sqrt{2}-5\approx2{,}1$.}
\end{ex}

\begin{ex}%[Dự án Tex hoá đề Tốt nghiệp 2025 + Phát triển 05 đề]%[Đỗ Đường Hiếu]%[0D0C2-9]
Điền $6$ số $1$, $2$, $3$, $4$, $5$, $6$ (mỗi số một lần) vào các ô tròn ở trên Hình $1$. Tính xác suất để điền các số đó vào các ô tròn để sao cho tổng các số ở mỗi cạnh của tam giác là bằng nhau (ví dụ ở Hình $2$, tổng các số ở mỗi cạnh đều bằng $10$). 
\begin{longtable}{m{0.3\textwidth}m{0.3\textwidth}}
\begin{tikzpicture}[scale=1,
>=stealth,auto,every node/.style={draw,circle}]
\node[minimum size=1cm] (a) at (0,0)  {};
\node[minimum size=1cm] (b) at (4,0)  {};
\node[minimum size=1cm] (c) at (60:4) {};
\node[minimum size=1cm] (m) at ($(a)!.5!(b)$) {};
\node[minimum size=1cm] (n) at ($(c)!.5!(b)$) {};
\node[minimum size=1cm] (p) at ($(a)!.5!(c)$) {};
\draw (a) edge (m);
\draw (m) edge (b);
\draw (b) edge (n);
\draw (n) edge (c);
\draw (c) edge (p);
\draw (p) edge (a);
\end{tikzpicture} &
\begin{tikzpicture}[scale=1,
>=stealth,auto,every node/.style={draw,circle}]
\node[minimum size=1cm] (a) at (0,0)  {$3$};
\node[minimum size=1cm] (b) at (4,0)  {$5$};
\node[minimum size=1cm] (c) at (60:4) {$1$};
\node[minimum size=1cm] (m) at ($(a)!.5!(b)$) {$2$};
\node[minimum size=1cm] (n) at ($(c)!.5!(b)$) {$4$};
\node[minimum size=1cm] (p) at ($(a)!.5!(c)$) {$6$};
\draw (a) edge (m);
\draw (m) edge (b);
\draw (b) edge (n);
\draw (n) edge (c);
\draw (c) edge (p);
\draw (p) edge (a);
\end{tikzpicture}\\
\centerline{Hình 1} & \centerline{Hình 2}
\end{longtable}
\shortans[oly]{$30$}
\loigiai{
\immini[thm]{Gọi các số điền vào là $A_1$, $A_2$, $A_3$,  $B_1$, $B_2$, $B_3$ như hình vẽ.\\
Ta có \begin{eqnarray*}
&&A_1+ B_2+ A_3= A_1+ B_3+ A_2= A_2+  B_1+ A_3\\
&\Leftrightarrow &\heva{& A_1+ B_2+ A_3= A_1+ B_3+ A_2 \\ &A_1+ B_2+ A_3= A_2+  B_1+ A_3}\\
&\Leftrightarrow & \heva{&  A_3- B_3= A_2- B_2 \\ &  A_1-  B_1= A_2- B_2}\\
&\Leftrightarrow &  A_1-  B_1= A_2- B_2= A_3- B_3.
\end{eqnarray*} }{\begin{tikzpicture}[scale=1,
>=stealth,auto,every node/.style={draw,circle}]
\node[minimum size=1cm] (a) at (0,0)  {$A_2$};
\node[minimum size=1cm] (b) at (4,0)  {$A_3$};
\node[minimum size=1cm] (c) at (60:4) {$A_1$};
\node[minimum size=1cm] (m) at ($(a)!.5!(b)$) {$B_1$};
\node[minimum size=1cm] (n) at ($(c)!.5!(b)$) {$B_2$};
\node[minimum size=1cm] (p) at ($(a)!.5!(c)$) {$B_3$};
\draw (a) edge (m);
\draw (m) edge (b);
\draw (b) edge (n);
\draw (n) edge (c);
\draw (c) edge (p);
\draw (p) edge (a);
\end{tikzpicture}}\noindent
Do $A_1$, $A_2$, $A_3$,  $B_1$, $B_2$, $B_3$ là một hoán vị của $1$, $2$, $3$, $4$, $5$, $6$.
Nên ta chỉ có các bộ sau thỏa mãn: $6-5 = 4-3 = 2-1$; $5-6 = 3-4 = 1-2$; $6-3 = 5-2 = 4-1$; $3-6 = 2-5 = 1-4$.\\
Ứng với mỗi bộ ở trên ta có $3!$ hoán vị các đỉnh $A_1$, $A_2$, $A_3$.
Và với mỗi cách chọn $A_1$, $A_2$, $A_3$ thì sẽ có duy nhất một cách chọn $B_1$, $B_2$, $B_3$.\\
Vậy có $3!\times 4=24$ cách điền thỏa mãn yêu cầu  tổng các số ở mỗi cạnh của tam giác là bằng nhau.\\
Khi đó $a=\dfrac{\mathrm{24}}{6!}=\dfrac{1}{30}\Rightarrow \dfrac{1}{a}=30$.
}
\end{ex}

\Closesolutionfile{ans}
=======
\begin{name}
	{\tenchude}
	{TOÁN 10}
	{LỚP TOÁN THẦY PHÁT}
	{I can accept failure, everyone fails at something. But I can’t accept not trying.}
\end{name}
\Opensolutionfile{ansbook}[ans/ansbookDe8]
\TN
\Opensolutionfile{ans}[ans/ansDe1-TN8]
\begin{ex}%[HSA-đợt 1,Huy Trần- Đăng Tạ]%[0D8N1-5]
Từ các chữ số $1$, $2$, $3$, $4$, $5$ có thể lập được bao nhiêu số tự nhiên có $3$ chữ số đôi một khác nhau?
\choice{\True $60$} {$125$} {$5$} {$10$}
\loigiai{
Lập số tự nhiên có $3$ chữ số đôi một khác nhau từ các chữ số $1$, $2$, $3$, $4$, $5$ có
\[A_3^5 = 60 ~\text {số}.\]
}
\end{ex}

\begin{ex}%[0D8N2-1]%[Dự án đề kiểm tra Toán khối 11 HKII NH23-24-Dot15-PhucChuong]%[THPT Nhu Van Lang - Hai Phong]
Công thức tính số chỉnh hợp chập $k$ của $n$ phần tử là:
\choice
{$A_n^k=\dfrac{n!}{\left(n-k\right)!k!}$}
{$C_n^k=\dfrac{n!}{\left(n-k\right)!}$}
{$C_n^k=\dfrac{n!}{\left(n-k\right)!k!}$}
{\True$A_n^k=\dfrac{n!}{\left(n-k\right)!}$}
\loigiai{
Công thức tính số chỉnh hợp chập $k$ của $n$ phần tử là $A_n^k=\dfrac{n!}{\left(n-k\right)!}$.
}
\end{ex}

\begin{ex}%[0D8N3-4]
Trong khai triển $\left(2a-b\right)^5$, hệ số của số hạng thứ $3$ bằng
\choice
{$-80$}
{\True $80$}
{$-10$}
{$10$}
\loigiai{
$
\begin{aligned}
\text{Ta có}
\left(2a-b\right)^5&=(2a)^5-5(2a)^4b+10(2a)^3b^2-10(2a)^2b^3+5(2a)b^4-b^5 \\
&=32a^5-80a^4b+80a^3b^2-40a^2b^3+10ab^4-b^5.
\end{aligned}
$
}
\end{ex}

\begin{ex}%[H]%[Nguyễn Văn Cường (Cường NV) - Dự án giảng K10 - K11]%[0D6H1-1]
Cho số $x = \dfrac{2}{7}$ và các giá trị gần đúng của $x$ là $0{,}28$; $0{,}29$; $0{,}286$; $0{,}3$. Hãy xác định sai số tuyệt đối trong từng trường hợp và cho biết giá trị gần đúng nào là tốt nhất.
\choice
{$0{,}28$}
{$0{,}29$}
{\True $0{,}286$}
{$0{,}3$}
\loigiai{
Ta có các sai số tuyệt đối là
\[\Delta_a=\left|\dfrac{2}{7}-0{,}28\right|=\dfrac{1}{175}; \Delta_b=\left|\dfrac{2}{7}-0{,}29\right|=\dfrac{3}{700} \Delta_c=\left|\dfrac{2}{7}-0{,}286\right|=\dfrac{1}{3500}; \Delta_d=\left|\dfrac{2}{7}-0{,}3\right|=\dfrac{1}{70}.\]
Vì $\Delta_c < \Delta_b < \Delta_a < \Delta_d$ nên $c=0{,}286$ là số gần đúng tốt nhất.}
\end{ex}

\begin{ex}%[0D6H4-2]
Nhiệt độ trung bình hàng tháng trong một năm được ghi lại trong bảng sau:
\begin{center}
\begin{tabular}{|c|c|c|c|c|c|c|c|c|c|c|c|c|}
\hline
Tháng &1&2&3&4&5&6&7&8&9&10&11&12\\
\hline
Nhiệt độ&16&20&24&28&30&30&29&25&25&20&18&16\\
\hline
\end{tabular}
\end{center}
Tìm khoảng tứ phân vị của bảng số liệu trên.
\choice
{$24,5$}
{$19$}
{$28,5$}
{\True $9,5$}
\loigiai{
Sắp xếp mẫu số liệu theo thứ tự không giảm:
\[16\hspace{1cm}16\hspace{1cm}18\hspace{1cm}20\hspace{1cm}20\hspace{1cm}24\hspace{1cm}25\hspace{1cm}25\hspace{1cm}28\hspace{1cm}29\hspace{1cm}30\hspace{1cm}30\]
Mẫu số liệu gồm $12$ giá trị nên trung vị là $Q_2=\dfrac{24+25}{2}=24,5$.\\
Nửa số liệu bên trái gồm $6$ giá trị. Do đó, $Q_1=\dfrac{18+20}{2}=19$.\\
Nửa số liệu bên phải gồm $6$ giá trị. Do đó, $Q_3=\dfrac{28+29}{2}=28,5$.\\
Vậy khoảng tứ phân vị của bảng số liệu trên là: $\Delta_Q=28,5-19=9,5$.
}
\end{ex}

\begin{ex}%[0D0N1-2]%[Dự án đề kiểm tra Toán 11 GHKII NH23-24- Hieu Hieu Minh Minh]%[THPT Toàn Thắng- Hải Phòng]
Không gian mẫu của phép thử \lq\lq Gieo đồng thời $2$ đồng xu cân đối và đồng chất\rq\rq\, là
\choice
{$\Omega=\left\{S, N, S, N\right\}$}
{\True $\Omega=\left\{SS, NN, SN, NS\right\}$}
{$\Omega=\left\{SS, NN\right\}$}
{$\Omega=\left\{S, N\right\}$}
\loigiai{
Gieo đồng thời $2$ đồng xu cân đối và đồng chất, ta được không gian mẫu là $\Omega=\left\{SS, NN, SN, NS\right\}$.
}
\end{ex}

\begin{ex}%[Dự án 17 - CK2 NH23-24 - TeamTeXHoa - ThinhNgo]%[0D0N1-3]
Gieo đồng tiền $2$ lần. Số phần tử của biến cố để mặt ngửa xuất hiện ít nhất $1$ lần là
\choice
{$6$}
{$5$}
{\True $3$}
{$4$}
\loigiai{
Không gian mẫu $\Omega=\{NN,NS,SN,SS\}$.\\
Gọi $A$ là biến cố mặt ngửa xuất hiện ít nhất $1$ lần. Khi đó $A=\{NN,NS,SN\}$ $\;\Rightarrow\; n(A)=3$.
}
\end{ex}

\begin{ex}%[0H9H3-2]
Trong mặt phẳng $Oxy$, phương trình tổng quát của đường thẳng đi qua hai điểm $A\left(3;-1\right)$ và $B\left(1;5\right)$ là
\choice
{$-x+3y+6=0$}
{$3x-y+6=0$}
{$3x-y+10=0$}
{\True $3x+y-8=0$}
\loigiai{
\begin{align*}
\heva{	& A\left(3;-1\right)\in AB \\
& {{\vec{u}}}_{AB}=\overrightarrow{AB}=\left(-2;6\right)\Rightarrow {{\vec{n}}}_{AB}=\left(3;1\right).}
\end{align*}
Phương trình đường thẳng $AB$ là \[3\left(x-3\right)+1\left(y+1\right)=0\Leftrightarrow 3x+y-8=0.\]
}
\end{ex}

\begin{ex}%[0H9H3-3]
Hai đường thẳng $\left(d_1\right) \colon \heva{&x=-2+5 t \\ &y=2 t}$ và $\left(d_2\right) \colon 4 x+3 y-18=0$ cắt nhau tại điểm có tọa độ
\choice
{\True $(2 ; 3)$}
{$(3 ; 2)$}
{$(1 ; 2)$}
{$(2 ; 1)$}
\loigiai{
Ta có
\[
\left(d_1\right) \colon \heva{&x=-2+5 t \\&y=2 t}
\Rightarrow
\left(d_1\right) \colon 2 x-5 y+4=0
\]
Gọi $ M $ là giao điểm của $ (d_1) $ và $ (d_2) $ nên tọa độ $ M $ là nghiệm của hệ phương trình
\[
\heva{&2 x-5 y+4=0 \\ &4 x+3 y-18=0}
\Leftrightarrow
\heva{&x=2 \\ &y=3.}
\]
Vậy $ M(2;3) $.
}
\end{ex}

\begin{ex}%[0H9H3-5]%[tex hóa đề ck2 - form 2025 - đợt 2- Duong Xuan Loi]
Khoảng cách từ điểm $M(3;-1)$ đến đường thẳng $\Delta\colon\heva{&x=-2+t\\&y=1+2t}$ nằm trong khoảng nào sau đây?
\choice
{$(1;3)$}
{$(3;5)$}
{$(7;9)$}
{\True $(5;7)$}
\loigiai{
Phương trình tổng quát đường thẳng $\Delta$ là $2x-y+5=0$.\\
Khoảng cách từ điểm. $M$. đến đường thẳng $\Delta$ là $\dfrac{\left|2\cdot 3-(-1)+5\right|}{\sqrt{2^2+(-1)^2}}=\dfrac{12\sqrt{5}}{5}\approx 5,4$.}
\end{ex}

\begin{ex}%[0H9N5-1]%[Dự án đề kiểm tra Toán Khối 10 CK2 NH23-24-Dot 16- Khắc Thiên]%[THPT Chuyên Hùng Vương - Phú Thọ]
Trong mặt phẳng $Oxy$, cho đường elip $(E)$ có phương trình  $\dfrac{x^2}{36}+\dfrac{y^2}{9}=1$. Tổng khoảng cách từ mỗi điểm trên elip tới hai tiêu điểm bằng
\choice
{$6$}
{$3$}
{$5$}
{\True $12$}
\loigiai{
Ta có phương trình của chính tắc elip  $(E)$ là $\dfrac{x^2}{36}+\dfrac{y^2}{9}=1$.\\
Suy ra $a=6$.\\
Khi đó, tổng khoảng cách từ mỗi điểm trên elip tới hai tiêu điểm bằng $2a=12$.
}
\end{ex}

\begin{ex}%[0H9N5-5]
Phương trình nào sau đây là phương trình chính tắc của hypebol?
\choice
{$x^2+\dfrac{y^2}{3^2}=1$}
{$\dfrac{x^2}{16}-y^2=-1$}
{$\dfrac{x^2}{25}-\dfrac{y^2}{9}=-1$}
{\True $x^2-\dfrac{y^2}{2}=1$}
\loigiai{
Phương trình
$x^2-\dfrac{y^2}{2}=1 \Leftrightarrow \dfrac{x^2}{1}-\dfrac{y^2}{2}=1$ là phương trình chính tắc của hypebol.
}
\end{ex}
\TNTF
\begin{ex}%[24-25 giảng K10 - K11, Phan Quốc Trí]%[0D8V2-6]
Cho đa giác đều $12$ cạnh. Xét tính đúng sai của các khẳng định sau.
\choiceTF
{\True Đa giác này có $54$ đường chéo}
{\True Có $220$ tam giác được tạo thành từ các đỉnh của đa giác}
{Số hình chữ nhật được tạo thành từ các đỉnh của đa giác là $165$}
{Từ các đỉnh của đa giác, có $120$ tam giác có đúng một cạnh là cạnh của đa giác}
\loigiai{
\begin{itemchoice}
\itemch Đúng. Số đường chéo của đa giác đều $12$ cạnh là:
$\mathrm{C}_{12}^2 - 12=54$ đường chéo.
\itemch Đúng. Số tam giác được tạo thành từ các đỉnh của đa giác là số cách chọn $3$ đỉnh từ $12$ đỉnh:
$\mathrm{C}_{12}^3=220$ tam giác.
\itemch Sai. Đa giác đều này có $6$ đường chéo đi qua tâm. Cứ hai đường chéo qua tâm cho ta một hình chữ nhật. Do đó có $\mathrm{C}_{6}^2 = 15$ hình chữ nhật tạo thành.
\itemch Sai. Trước tiên ta chọn $1$ cạnh trong $12$ cạnh của đa giác nên có $\mathrm{C}_{12}^1$ cách chọn; tiếp theo chọn $1$ đỉnh còn lại trong $8$ đỉnh (trừ $2$ đỉnh tạo nên cạnh đã chọn và $2$ đỉnh liền kề với cạnh đã chọn) nên có $\mathrm{C}_{8}^1$ các chọn. Do đó Số tam giác có một cạnh là cạnh của đa giác là:
$\mathrm{C}_{12}^1\times\mathrm{C}_{8}^1=96$ tam giác.
\end{itemchoice}
}
\end{ex}

\begin{ex}%[0H9H3-2]%[Dự án đề 4 phần]$%[Tex hóa: Lê Thị Thúy Hằng]
Trong mặt phẳng toạ độ $Oxy$, cho hai điểm $A(-2;2)$, $B(3;4)$.
\choiceTF
{Đường thẳng $AB$ có vectơ chỉ phương là $\overrightarrow{AB}=(2;5)$}
{Trung điểm của đoạn thẳng $AB$ có tọa độ là $(1;6)$}
{\True Phương trình tổng quát của đường thẳng $AB$ là $2x-5y+14=0$}
{Phương trình tham số của đường thẳng đi qua $M(-1;1)$ và song song với $AB$ là $\heva{&x=-1+2t\\&y=1+5t}$}
\loigiai{
\begin{itemchoice}
\itemch {\bf Sai}.\\
Đường thẳng $AB$ có vectơ chỉ phương là $\overrightarrow{AB}=(5;2)$.
\itemch {\bf Sai}.\\
Trung điểm của đoạn thẳng $AB$ có tọa độ là $\left(\dfrac{1}{2};3\right)$.
\itemch {\bf Đúng}.\\
Phương trình tổng quát của đường thẳng $AB$ đi qua $A(-2;2)$ và có vectơ pháp tuyến $\overrightarrow{n}=(2;-5)$ là \[2(x+2)-5(y-2)=0 \Leftrightarrow 2x-5y+14=0.\]
\itemch {\bf Sai}.\\
Đường thẳng đi qua $M(-1;1)$, song song với đường thẳng $AB$ nên nhận $\overrightarrow{AB}=(5;2)$ làm vectơ chỉ phương, có phương trình tham số là $\heva{&x=-1+5t\\&y=1+2t.}$
\end{itemchoice}
}
\end{ex}
\TNSA
\begin{ex}%[0D8V1-2]%[Dự án đề kiểm tra Toán 10 CHKII NH23-24- Ngô Tất Thành]%[THPT Chuyên Hùng Vương - Tỉnh Phú Thọ]
Đội thanh niên xung kích của một trường trung học phổ thông có $12$ học sinh trong đó có $9$ học sinh nam và $3$ học sinh nữ. Đoàn trường cần chọn một nhóm $5$ học sinh đi làm nhiệm vụ sao cho phải có $1$ đội trưởng nam, $1$ đội phó nam và có ít nhất $1$ nữ. Hỏi có bao nhiêu cách?
\shortans{$6120$}
\loigiai{
Chọn $1$ nam làm đội trưởng có $\mathrm{C}_9^1$ cách.\\
Chọn $1$ nam làm đội phó trong 8 nam còn lại có $\mathrm{C}_8^1$ cách.\\
Chọn $3$ học sinh sao cho có ít nhất $1$ nữ có $\mathrm{C}_{10}^3-\mathrm{C}_7^3=85$.\\
Số cách chọn $5$ bạn thỏa yêu cầu bài toán là $\mathrm{C}_9^1 \cdot \mathrm{C}_8^1 \cdot 85=6120$ cách.
}
\end{ex}

\begin{ex}%[0D8N3-4]
Tổng các hệ số của các số hạng trong khai triển nhị thức $(x-2y)^{2\,020}$ bằng
\choice{$2\,021$}{$2\,020$}{$-1$}{\True $1$}
\loigiai{
Ta có $(x-2y)^{2\,020}=\sum\limits_{k=0}^{2\,020}\mathrm{C}_{2\,020}^kx^{2\,020-k}(-2y)^k=\sum\limits_{k=0}^{2\,020} a_kx^{2\,020-k}y^k$.\\
Chọn $x=y=1$ ta có $\sum\limits_{k=0}^{2\,020} a_k=(-1)^{2\,020}=1$.
}
\end{ex}

\begin{ex}%[De-chuan-hoa-so-8]%[Duong Xuan Loi]%[0D6H3-2]
Bảng dưới đây thống kê nhiệt độ (đơn vị: ${}^{\circ}{C}$) ở thành phố Hồ Chí Minh ngày $03 / 06 / 2021$ sau một số lần đo
\begin{center}
\begin{tabular}{|l|c|c|c|c|c|c|c|c|}
\hline Giờ đo & $1 $ h
& $4 $ h & $7 $ h & $10 $ h & $13 $ h & $16 $ h & $19 $ h & $22 $ h \\
\hline Nhiệt độ $\left({}^{\circ}{C}\right)$ & $27$ & $26$ & $28$ & $32$ & $34$ & $35$ & $30$ & $28$ \\
\hline
\end{tabular}
\end{center}
Tìm số trung bình của mẫu số liệu (làm tròn kết quả đến hàng phần trăm).
\shortans{$29$}
\loigiai{
Sắp xếp dãy số liệu nhiệt độ theo thứ tự không giảm ta được: $26;27;28;28;30;32;34;35$.\\
Nhiệt độ trung bình $\overline{x}=\dfrac{26+27+28\cdot 2+30+32+34+35}{8}=30$.\\
Số trung vị $M_e=\dfrac{28+30}{2}=29$.
}
\end{ex}

\begin{ex}%[0H9V4-2]
Trong mặt phẳng tọa độ $Oxy$, cho điểm $A(-1;1)$, $B(0;-2)$, $C(0;2)$. Đường tròn ngoại tiếp tam giác $ ABC $ có phương trình là $ (x-a)^2+(y-b)^2=R^2 $. Tính $ S=a+b+R^2 $.
\shortans{$6$}
\loigiai{
Giả sử tâm đường tròn là điểm $I(a;b)$.\\
Ta có \allowdisplaybreaks
\begin{eqnarray*}
& &IA=IB=IC\Leftrightarrow IA^2=IB^2=IC^2\\
&\Leftrightarrow&\heva{& (-1-a)^2+(1-b)^2=(0-a)^2+(-2-b)^2\\ & (0-a)^2+(-2-b)^2=(0-a)^2+(2-b)^2}\\
&\Leftrightarrow& \heva{& a^2+b^2+2a-2b+2=a^2+b^2+4b+4\\ & a^2+b^2+4b+4=a^2+b^2-4b+4}\\
&\Leftrightarrow& \heva{& 2a-2b=4b+2\\ & b=0}\Leftrightarrow\heva{& a=1\\ & b=0.}
\end{eqnarray*}
Đường tròn tâm $I(1;0)$, bán kính $R=IC=\sqrt{a^2+b^2-4b+4}=\sqrt{5}$.\\
Suy ra phương trình đường tròn đi qua ba điểm $A$, $B$, $C$ có dạng $(x-1)^2+y^2=5$.\\
Vậy $a=1$, $b=0$, $R^2=5$, $a+b+R^2=6$.
}
\end{ex}
\TL
\begin{ex}%[0H9H3-2]%[Dự án đề kiểm tra Toán 10 HK2 NH23-24-BCTuan]%[THPT Nguyễn Thị Minh Khai-TPHCM]
Trong mặt phẳng $Oxy$, cho $A(1;2)$, $B(2;3)$, $C(6;6)$. Viết phương trình tổng quát của đường thẳng $\Delta$ đi qua điểm $A$ và song song với $BC$.
\loigiai{
Ta có $\overrightarrow{BC}=(4;3)$ suy ra $\overrightarrow{n}=(3;-4)$ là véc-tơ pháp tuyến của $BC$.\\
Vì $\Delta$ song song với $BC$ nên $\overrightarrow{n}=(3;-4)$ cũng là véc-tơ pháp tuyến của $\Delta$.\\
Ta có phương trình tổng quát của $\Delta$ là $3(x-1)-4(y-2)=0\Leftrightarrow 3x-4y+5=0$.
}
\end{ex}

\begin{ex}%[0H9V4-5]%[Dự án đề kiểm tra Toán 11 GHKI NH23-24- Phan Anh]%[THPT - Tp HCM]
Trong mặt phẳng tọa độ $Oxy$, cho điểm $I(2;-1)$ và đường thẳng $d\colon3x+y+5=0$. Đường tròn $(C)$ có tâm $I$ và đường kính bằng $4\sqrt{5}$. Đường thẳng $\Delta\colon x+by+c=0$, $c>0$ vuông góc với đường thẳng $d$ và cắt $(C)$ tại hai điểm phân biệt $A$, $B$ sao cho tam giác $IAB$ tù và có diện tích bằng $5\sqrt{3}$. Giá trị của $c$ bằng bao nhiêu? (Làm tròn kết quả đến hàng phần chục)
\loigiai{\immini{Đường tròn $(C)$ có đường kính bằng $4\sqrt{5}$ nên có bán kính là $R=2\sqrt{5}$.\\
Đường thẳng $\Delta$ vuông góc với đường thẳng $d$ nên $\Delta\colon x-3y+c=0$.\\
Tam giác $IAB$ có diện tích là $5\sqrt{3}$ nên
\[S=\dfrac{1}{2}\cdot IA\cdot IB\cdot\sin\widehat{AIB}\Leftrightarrow5\sqrt{3}=\dfrac{R^2\sin\widehat{AIB}}{2}\Leftrightarrow\sin\widehat{AIB}=\dfrac{\sqrt{3}}{2}.\]}
{\begin{tikzpicture}[>=stealth,line join=round,line cap=round,font=\footnotesize,scale=0.3]
\path (0,0) coordinate (I)
($(I)+({5*cos(30)},{5*sin(30)})$) coordinate(A)
($(I)!1!120:(A)$) coordinate(B)
($(A)!0.5!(B)$) coordinate(H)
;
\draw (I)circle(5);
\draw (H)--(I)--(B)--(A)--(I);
\foreach \p / \r in {A/45,I/-90,B/135,H/90}
\fill (\p) circle (1.2pt) node[shift={(\r:3mm)}]{$\p$};
\end{tikzpicture}}
\noindent
Do $\triangle IAB$ tù nên $\widehat{AIB}=120^\circ$, suy ra khoảng cách từ $I$ đến $\Delta$ là $IH=\dfrac{R}{2}=\sqrt{5}$.\\
Từ đó ta có $IH=\mathrm{d}(I,\Delta)=\dfrac{|2+3+c|}{\sqrt{1+9}}=\sqrt{5}\Leftrightarrow\hoac{&c=5\sqrt{2}-5\\&c=-5\sqrt{2}-5.}$\\
Do $c$ dương nên nhận $c=5\sqrt{2}-5\approx2{,}1$.}
\end{ex}

\begin{ex}%[Dự án Tex hoá đề Tốt nghiệp 2025 + Phát triển 05 đề]%[Đỗ Đường Hiếu]%[0D0C2-9]
Điền $6$ số $1$, $2$, $3$, $4$, $5$, $6$ (mỗi số một lần) vào các ô tròn ở trên Hình $1$. Tính xác suất để điền các số đó vào các ô tròn để sao cho tổng các số ở mỗi cạnh của tam giác là bằng nhau (ví dụ ở Hình $2$, tổng các số ở mỗi cạnh đều bằng $10$). 
\begin{longtable}{m{0.3\textwidth}m{0.3\textwidth}}
\begin{tikzpicture}[scale=1,
>=stealth,auto,every node/.style={draw,circle}]
\node[minimum size=1cm] (a) at (0,0)  {};
\node[minimum size=1cm] (b) at (4,0)  {};
\node[minimum size=1cm] (c) at (60:4) {};
\node[minimum size=1cm] (m) at ($(a)!.5!(b)$) {};
\node[minimum size=1cm] (n) at ($(c)!.5!(b)$) {};
\node[minimum size=1cm] (p) at ($(a)!.5!(c)$) {};
\draw (a) edge (m);
\draw (m) edge (b);
\draw (b) edge (n);
\draw (n) edge (c);
\draw (c) edge (p);
\draw (p) edge (a);
\end{tikzpicture} &
\begin{tikzpicture}[scale=1,
>=stealth,auto,every node/.style={draw,circle}]
\node[minimum size=1cm] (a) at (0,0)  {$3$};
\node[minimum size=1cm] (b) at (4,0)  {$5$};
\node[minimum size=1cm] (c) at (60:4) {$1$};
\node[minimum size=1cm] (m) at ($(a)!.5!(b)$) {$2$};
\node[minimum size=1cm] (n) at ($(c)!.5!(b)$) {$4$};
\node[minimum size=1cm] (p) at ($(a)!.5!(c)$) {$6$};
\draw (a) edge (m);
\draw (m) edge (b);
\draw (b) edge (n);
\draw (n) edge (c);
\draw (c) edge (p);
\draw (p) edge (a);
\end{tikzpicture}\\
\centerline{Hình 1} & \centerline{Hình 2}
\end{longtable}
\shortans[oly]{$30$}
\loigiai{
\immini{Gọi các số điền vào là $A_1$, $A_2$, $A_3$,  $B_1$, $B_2$, $B_3$ như hình vẽ.\\
Ta có \begin{eqnarray*}
&&A_1+ B_2+ A_3= A_1+ B_3+ A_2= A_2+  B_1+ A_3\\
&\Leftrightarrow &\heva{& A_1+ B_2+ A_3= A_1+ B_3+ A_2 \\ &A_1+ B_2+ A_3= A_2+  B_1+ A_3}\\
&\Leftrightarrow & \heva{&  A_3- B_3= A_2- B_2 \\ &  A_1-  B_1= A_2- B_2}\\
&\Leftrightarrow &  A_1-  B_1= A_2- B_2= A_3- B_3.
\end{eqnarray*} }{\begin{tikzpicture}[scale=1,
>=stealth,auto,every node/.style={draw,circle}]
\node[minimum size=1cm] (a) at (0,0)  {$A_2$};
\node[minimum size=1cm] (b) at (4,0)  {$A_3$};
\node[minimum size=1cm] (c) at (60:4) {$A_1$};
\node[minimum size=1cm] (m) at ($(a)!.5!(b)$) {$B_1$};
\node[minimum size=1cm] (n) at ($(c)!.5!(b)$) {$B_2$};
\node[minimum size=1cm] (p) at ($(a)!.5!(c)$) {$B_3$};
\draw (a) edge (m);
\draw (m) edge (b);
\draw (b) edge (n);
\draw (n) edge (c);
\draw (c) edge (p);
\draw (p) edge (a);
\end{tikzpicture}}\noindent
Do $A_1$, $A_2$, $A_3$,  $B_1$, $B_2$, $B_3$ là một hoán vị của $1$, $2$, $3$, $4$, $5$, $6$.
Nên ta chỉ có các bộ sau thỏa mãn: $6-5 = 4-3 = 2-1$; $5-6 = 3-4 = 1-2$; $6-3 = 5-2 = 4-1$; $3-6 = 2-5 = 1-4$.\\
Ứng với mỗi bộ ở trên ta có $3!$ hoán vị các đỉnh $A_1$, $A_2$, $A_3$.
Và với mỗi cách chọn $A_1$, $A_2$, $A_3$ thì sẽ có duy nhất một cách chọn $B_1$, $B_2$, $B_3$.\\
Vậy có $3!\times 4=24$ cách điền thỏa mãn yêu cầu  tổng các số ở mỗi cạnh của tam giác là bằng nhau.\\
Khi đó $a=\dfrac{\mathrm{24}}{6!}=\dfrac{1}{30}\Rightarrow \dfrac{1}{a}=30$.
}
\end{ex}

\Closesolutionfile{ans}
>>>>>>> 33890a9ce1cfbd3079ed36a93e836cf3ef1fde30
\Closesolutionfile{ansbook}